\section{Обзор раздела}

Данный раздел описывает практические аспекты реализации архитектуры DREAM-Net. Цель — предоставить разработчикам достаточную информацию для корректной реализации без двусмысленностей, сохраняя при этом гибкость для оптимизации под конкретные платформы.

\textbf{Важное замечание:} Все представленные значения гиперпараметров являются предварительными рекомендациями на основе теоретического анализа и ограниченных предварительных тестов. Они могут требовать настройки под конкретные задачи и датасеты.

\section{Программный интерфейс (API Specification)}

\subsection{Основной класс модели}

\begin{lstlisting}[language=Python, caption={Интерфейс класса DREAMNet}, label={lst:dream_class}]
class DREAMNet:
    """
    DREAM-Net: Dynamic Recall and Elastic Adaptive Memory

    Архитектура с онлайн-пластичностью и двухмасштабной памятью.
    """

    def __init__(self, config: DREAMConfig):
        """
        Инициализация модели.

        Args:
            config: Конфигурация с гиперпараметрами
        """
        pass

    def step(self, x_t: Tensor, state: Optional[State] = None) -> StepOutput:
        """
        Один шаг инференса с возможностью адаптации.

        Args:
            x_t: Входной сигнал на шаге t, shape (d_in,)
            state: Скрытое состояние (опционально, создаётся если None)

        Returns:
            StepOutput с полями:
                - h_next: Новое скрытое состояние
                - prediction: Предсказание модели
                - surprise: Коэффициент удивления S_t
                - mode: Текущий режим (Routine/Learning/Shock)
        """
        pass

    def sleep(self, priority: int = 0) -> SleepResult:
        """
        Асинхронная фаза консолидации памяти.

        Args:
            priority: Приоритет задачи (для планировщика)

        Returns:
            SleepResult с полями:
                - consolidated: Количество консолидированных паттернов
                - duration: Длительность фазы сна (мс)
                - success: Флаг успешного завершения
        """
        pass

    def get_state(self) -> ModelState:
        """
        Получение текущего состояния для сериализации.

        Returns:
            ModelState с полями:
                - slow_weights: W (медленные веса)
                - fast_weights: U (быстрые веса)
                - context: c_t (контекстный вектор)
                - statistics: μ, σ (статистика ошибки)
                - config: Конфигурация модели
        """
        pass

    def load_state(self, state: ModelState) -> bool:
        """
        Загрузка состояния из сериализованного представления.

        Args:
            state: ModelState для загрузки

        Returns:
            True если загрузка успешна, False иначе
        """
        pass

    def reset(self, hard: bool = False):
        """
        Сброс состояния модели.

        Args:
            hard: Если True, сбрасывает также быстрые веса и статистику
        """
        pass
\end{lstlisting}

\subsection{Структуры данных}

\begin{lstlisting}[language=Python, caption={Структуры данных}, label={lst:data_structures}]
@dataclass
class DREAMConfig:
    """Конфигурация модели"""
    d_state: int = 256          # Размерность скрытого состояния
    d_ctx: int = 512            # Размерность контекста (степень двойки)
    d_in: int = 80              # Размерность входа (например, mel-спектр)
    rank: int = 32              # Ранг быстрых весов
    n_layers: int = 4           # Количество слоёв иерархии

    # Surprise Gate
    tau_base: float = 0.5       # Базовый порог удивления
    gamma: float = 0.5          # Температура сигмоиды
    beta_stat: float = 0.05     # Сглаживание статистики

    # State Machine
    theta_learn_enter: float = 0.7    # Вход в Learning (в сигмах)
    theta_learn_exit: float = 0.3     # Выход из Learning
    theta_shock_enter: float = 2.5    # Вход в Shock
    theta_shock_exit: float = 1.8     # Выход из Shock
    shock_window_ms: int = 8          # Окно подтверждения шока (мс)

    # Plasticity
    eta_plasticity: float = 0.01      # Скорость пластичности
    lambda_decay: float = 0.001       # Затухание быстрых весов
    u_max: float = 10.0               # Максимальная норма U

    # Sleep
    sleep_trigger_events: int = 100   # Событий для триггера сна
    sleep_trigger_time_sec: int = 300 # Время для триггера сна
    zeta_sleep: float = 0.1           # Скорость консолидации
    lambda_ewc: float = 0.5           # Коэффициент EWC

    # Safety
    ramp_duration_ms: int = 50        # Длительность рампы восстановления
    shock_max_duration_ms: int = 500  # Максимальная длительность шока
    sigma_min: float = 1e-4           # Минимальная дисперсия


@dataclass
class State:
    """Скрытое состояние модели"""
    h: List[Tensor]           # Скрытые состояния по слоям
    U: List[Tensor]           # Быстрые веса по слоям
    c: Tensor                 # Контекстный вектор
    mu: float                 # Бегущее среднее ошибки
    sigma: float              # Бегущее отклонение ошибки
    mode: int                 # Текущий режим (0=Routine, 1=Learning, 2=Shock)
    alpha_ramp: float         # Коэффициент рампы
    shock_counter: int        # Счётчик окна шока
    tau_base: float           # Адаптивный базовый порог


@dataclass
class StepOutput:
    """Результат шага инференса"""
    h_next: List[Tensor]
    prediction: Tensor
    surprise: float
    mode: int
    error_norm: float


@dataclass
class SleepResult:
    """Результат фазы сна"""
    consolidated: int
    duration_ms: float
    success: bool
    weights_updated: bool
\end{lstlisting}

\subsection{Callback-интерфейсы}

\begin{lstlisting}[language=Python, caption={Callback-интерфейсы}, label={lst:callbacks}]
class DREAMCallback:
    """Интерфейс для событий модели"""

    def on_mode_change(self, old_mode: int, new_mode: int, step: int):
        """Вызывается при смене режима работы"""
        pass

    def on_sleep_start(self, trigger_reason: str):
        """Вызывается при начале фазы сна"""
        pass

    def on_sleep_end(self, result: SleepResult):
        """Вызывается при завершении фазы сна"""
        pass

    def on_anomaly_detected(self, surprise: float, step: int):
        """Вызывается при детекции аномалии (S_t > threshold)"""
        pass

    def on_numerical_warning(self, tensor_name: str, issue: str):
        """Вызывается при детекции численных проблем (NaN, Inf, clipping)"""
        pass
\end{lstlisting}

\section{Гиперпараметры по умолчанию}

\subsection{Архитектурные параметры}

\begin{table}[h]
\centering
\caption{Архитектурные гиперпараметры}
\label{tab:hyper_arch}
\begin{tabular}{lcp{2cm}p{2.5cm}p{4cm}}
\toprule
\textbf{Параметр} & \textbf{Символ} & \textbf{Значение} & \textbf{Диапазон} & \textbf{Обоснование} \\
\midrule
Размерность состояния & $d_{state}$ & 256 & 128–512 & Баланс ёмкости и эффективности \\
Размерность контекста & $d_{ctx}$ & 512 & 256–1024 & Степень двойки для FFT \\
Ранг быстрых весов & $r$ & 32 & 16–64 & Низкоранговая аппроксимация \\
Количество слоёв & $L$ & 4 & 3–6 & Иерархия предиктивного кодирования \\
Размерность входа & $d_{in}$ & 80 & --- & Зависит от фич (mel-спектр) \\
\bottomrule
\end{tabular}
\end{table}

\subsection{Параметры Surprise Gate}

\begin{table}[h]
\centering
\caption{Параметры Surprise Gate}
\label{tab:hyper_surprise}
\begin{tabular}{lcp{2cm}p{2.5cm}p{4cm}}
\toprule
\textbf{Параметр} & \textbf{Символ} & \textbf{Значение} & \textbf{Диапазон} & \textbf{Обоснование} \\
\midrule
Базовый порог & $\tau_{base}$ & 0.5 & 0.3–0.7 & Средняя чувствительность \\
Температура & $\gamma$ & 0.5 & 0.3–1.0 & Крутизна сигмоиды \\
Сглаживание статистики & $\beta$ & 0.05 & 0.01–0.1 & Скорость адаптации $\mu, \sigma$ \\
Минимальная дисперсия & $\sigma_{min}$ & $10^{-4}$ & $10^{-5}$–$10^{-3}$ & Защита от деления на ноль \\
\bottomrule
\end{tabular}
\end{table}

\subsection{Параметры State Machine}

\begin{table}[h]
\centering
\caption{Параметры машины состояний}
\label{tab:hyper_fsm}
\begin{tabular}{lcp{2cm}p{2.5cm}p{4cm}}
\toprule
\textbf{Параметр} & \textbf{Символ} & \textbf{Значение} & \textbf{Диапазон} & \textbf{Обоснование} \\
\midrule
Вход в Learning & $\theta_{learn}^{enter}$ & $0.7\sigma$ & 0.5–1.0$\sigma$ & Умеренная новизна \\
Выход из Learning & $\theta_{learn}^{exit}$ & $0.3\sigma$ & 0.2–0.5$\sigma$ & Гистерезис 0.4$\sigma$ \\
Вход в Shock & $\theta_{shock}^{enter}$ & $2.5\sigma$ & 2.0–3.0$\sigma$ & Аномалия \\
Выход из Shock & $\theta_{shock}^{exit}$ & $1.8\sigma$ & 1.5–2.0$\sigma$ & Гистерезис 0.7$\sigma$ \\
Окно подтверждения & $T_{win}$ & 8 мс & 5–15 мс & Фильтрация импульсов \\
Макс. длительность шока & $T_{shock\_max}$ & 500 мс & 200–1000 мс & Защита от зависания \\
\bottomrule
\end{tabular}
\end{table}

\subsection{Параметры пластичности}

\begin{table}[h]
\centering
\caption{Параметры пластичности}
\label{tab:hyper_plasticity}
\begin{tabular}{lcp{2cm}p{2.5cm}p{4cm}}
\toprule
\textbf{Параметр} & \textbf{Символ} & \textbf{Значение} & \textbf{Диапазон} & \textbf{Обоснование} \\
\midrule
Скорость пластичности & $\eta$ & 0.01 & 0.001–0.1 & Скорость обновления $U$ \\
Затухание быстрых весов & $\lambda_{decay}$ & 0.001 & 0.0001–0.01 & Возврат к $U_{target}$ \\
Макс. норма $U$ & $U_{max}$ & 10.0 & 5.0–20.0 & Ограничение весов \\
Скорость консолидации & $\zeta_{sleep}$ & 0.1 & 0.05–0.2 & Перенос $U \to W$ \\
Коэффициент EWC & $\lambda_{EWC}$ & 0.5 & 0.1–1.0 & Защита от забывания \\
\bottomrule
\end{tabular}
\end{table}

\subsection{Параметры восстановления}

\begin{table}[h]
\centering
\caption{Параметры восстановления}
\label{tab:hyper_recovery}
\begin{tabular}{lcp{2cm}p{2.5cm}p{4cm}}
\toprule
\textbf{Параметр} & \textbf{Символ} & \textbf{Значение} & \textbf{Диапазон} & \textbf{Обоснование} \\
\midrule
Длительность рампы & $T_{ramp}$ & 50 мс & 20–100 мс & Плавное восстановление \\
Затухание порога & $\delta_{decay}$ & 0.01 & 0.005–0.02 & Возврат $\tau_{base}$ к норме \\
Повышение порога & $\delta_{up}$ & 0.10 & 0.05–0.15 & Адаптация в шоке \\
Потолок порога & $\tau_{max}$ & $3.0 \cdot \tau_0$ & 2–5$\cdot \tau_0$ & Макс. адаптация \\
Пол порога & $\tau_{min}$ & $0.5 \cdot \tau_0$ & 0.3–0.7$\cdot \tau_0$ & Мин. адаптация \\
\bottomrule
\end{tabular}
\end{table}

\section{Структура кода (Recommended Code Structure)}

\begin{lstlisting}[language=bash, caption={Рекомендуемая структура проекта}, label={lst:project_structure}]
dream_net/
├── __init__.py
├── config.py                 # DREAMConfig, presets
├── model/
│   ├── __init__.py
│   ├── dream_net.py          # Основной класс DREAMNet
│   ├── layers/
│   │   ├── __init__.py
│   │   ├── predictive_layer.py    # Слой предиктивного кодирования
│   │   ├── surprise_gate.py       # Surprise Gate
│   │   └── context_module.py      # FHRR контекст
│   ├── memory/
│   │   ├── __init__.py
│   │   ├── fast_weights.py        # Быстрые веса U
│   │   ├── slow_weights.py        # Медленные веса W
│   │   └── double_buffer.py       # Double Buffering
│   └── state_machine/
│       ├── __init__.py
│       ├── fsm.py                 # Конечный автомат
│       └── modes.py               # Режимы (Routine/Learning/Shock)
├── training/
│   ├── __init__.py
│   ├── sleep_cycle.py             # Фаза сна
│   ├── consolidation.py           # Консолидация
│   └── ewc.py                     # EWC-регуляризация
├── utils/
│   ├── __init__.py
│   ├── statistics.py              # Robust Sigma
│   ├── numerical.py               # Clipping, stability checks
│   └── callbacks.py               # Callback-интерфейсы
├── tests/
│   ├── __init__.py
│   ├── test_step.py               # Тесты шага инференса
│   ├── test_sleep.py              # Тесты фазы сна
│   ├── test_stability.py          # Тесты численной устойчивости
│   └── test_concurrency.py        # Тесты потокобезопасности
└── examples/
    ├── __init__.py
    ├── basic_asr.py               # Пример ASR
    ├── speaker_adaptation.py      # Пример адаптации к диктору
    └── noise_robustness.py        # Пример работы с шумом
\end{lstlisting}

\section{Зависимости и требования}

\subsection{Программные зависимости}

\begin{table}[h]
\centering
\caption{Программные зависимости}
\label{tab:dependencies}
\begin{tabular}{lp{3cm}p{6cm}}
\toprule
\textbf{Зависимость} & \textbf{Минимальная версия} & \textbf{Назначение} \\
\midrule
Python & 3.9 & Язык реализации \\
PyTorch & 2.0 & Тензорные операции, FFT \\
NumPy & 1.20 & Вспомогательные вычисления \\
threading/asyncio & --- & Асинхронный сон \\
\bottomrule
\end{tabular}
\end{table}

\subsection{Аппаратные требования}

\begin{table}[h]
\centering
\caption{Аппаратные требования}
\label{tab:hardware}
\begin{tabular}{lpp{6cm}}
\toprule
\textbf{Компонент} & \textbf{Минимум} & \textbf{Рекомендуется} \\
\midrule
CPU & 4 ядра & 8+ ядер (для асинхронного сна) \\
RAM & 2 GB & 4+ GB (Double Buffering) \\
GPU & Опционально & CUDA 11+ для ускорения FFT \\
Storage & 100 MB & 500 MB (для логов и чекпоинтов) \\
\bottomrule
\end{tabular}
\end{table}

\subsection{Оценка вычислительной сложности}

\begin{table}[h]
\centering
\caption{Вычислительная сложность операций}
\label{tab:complexity}
\begin{tabular}{lp{5cm}p{4cm}}
\toprule
\textbf{Операция} & \textbf{Сложность} & \textbf{Примечание} \\
\midrule
Шаг инференса & $O(L \cdot d_{state} \cdot r)$ & Линейная от ранга \\
Обновление контекста & $O(d_{ctx} \log d_{ctx})$ & FFT/IFFT \\
Фаза сна & $O(|W| + |U|)$ & Только при триггере \\
Память на шаг & $O(L \cdot d_{state} + d_{ctx})$ & Константна от длины последовательности \\
\bottomrule
\end{tabular}
\end{table}

\section{Интеграционные рекомендации}

\subsection{Интеграция с ASR-пайплайном}

\begin{lstlisting}[language=Python, caption={Пример интеграции с аудио-пайплайном}, label={lst:asr_integration}]
# Пример интеграции с аудио-пайплайном
import torch
from dream_net import DREAMNet, DREAMConfig

# Инициализация
config = DREAMConfig(d_in=80)  # 80 mel-фильтров
model = DREAMNet(config)

# Обработка аудио-потока
for audio_frame in audio_stream:
    # Предобработка (mel-спектр)
    x_t = extract_mel_features(audio_frame)

    # Шаг инференса
    output = model.step(x_t)

    # Использование предсказания
    if output.surprise > 0.8:
        logger.warning(f"High surprise detected: {output.surprise}")

    # Декодирование выхода (интеграция с CTC/Transducer)
    phoneme_probs = decoder(output.prediction)

    # Асинхронный сон (в отдельном потоке)
    if sleep_trigger.check():
        threading.Thread(target=model.sleep).start()
\end{lstlisting}

\subsection{Серверная развёртка}

Для серверной развёртки рекомендуется:

\begin{enumerate}
    \item \textbf{Контейнеризация:} Docker с предустановленными зависимостями.
    \item \textbf{Мониторинг:} Экспорт метрик (surprise\_rate, mode\_distribution, sleep\_frequency) в Prometheus.
    \item \textbf{Логирование:} Структурированные логи (JSON) с возможностью фильтрации по режимам.
    \item \textbf{Graceful shutdown:} Корректное завершение фазы сна перед остановкой сервиса.
\end{enumerate}

\subsection{Edge-развёртка}

Для устройств с ограниченными ресурсами:

\begin{enumerate}
    \item \textbf{Квантование:} INT8 для медленных весов $W$ (быстрые веса $U$ остаются FP32).
    \item \textbf{Отключение отладки:} Проверки \texttt{isfinite()} только в debug-сборках.
    \item \textbf{Упрощённый сон:} Синхронная фаза сна с блокировкой инференса.
    \item \textbf{Снижение ранга:} $r = 16$ вместо 32 для экономии памяти.
\end{enumerate}

\section{Отладка и диагностика}

\subsection{Рекомендуемые метрики для логирования}

\begin{table}[h]
\centering
\caption{Метрики для логирования}
\label{tab:logging_metrics}
\begin{tabular}{lp{4cm}p{5cm}}
\toprule
\textbf{Метрика} & \textbf{Частота} & \textbf{Назначение} \\
\midrule
\texttt{surprise\_mean} & Каждые 100 шагов & Средняя новизна \\
\texttt{mode\_distribution} & Каждые 1000 шагов & Доля времени в режимах \\
\texttt{U\_norm\_mean} & Каждые 100 шагов & Дрифт быстрых весов \\
\texttt{sleep\_frequency} & Каждые 10 минут & Частота консолидации \\
\texttt{clipping\_count} & Каждые 100 шагов & Срабатывания защит \\
\texttt{nan\_inf\_count} & Каждый шаг & Численные ошибки \\
\bottomrule
\end{tabular}
\end{table}

\subsection{Чек-лист перед продакшном}

\begin{itemize}
    \item[$\square$] Все unit-тесты проходят (test\_step, test\_sleep, test\_stability)
    \item[$\square$] Интеграционные тесты с реальными данными
    \item[$\square$] Стресс-тест (100+ часов непрерывной работы без NaN/Inf)
    \item[$\square$] Проверка Double Buffering на гонки данных
    \item[$\square$] Валидация на независимом тестовом наборе
    \item[$\square$] Документация API актуализирована
    \item[$\square$] Конфигурационные файлы вынесены из кода
\end{itemize}

\section{Известные ограничения и будущая работа}

\subsection{Текущие ограничения}

\begin{table}[h]
\centering
\caption{Текущие ограничения}
\label{tab:limitations}
\begin{tabular}{lp{5cm}p{4cm}}
\toprule
\textbf{Ограничение} & \textbf{Влияние} & \textbf{Временное решение} \\
\midrule
FFT на CPU медленно & Задержка на Edge & Использовать rfft, снизить $d_{ctx}$ \\
Сон блокирует на Edge & Фризы инференса & Синхронный сон с приоритетом \\
Гиперпараметры требуют настройки & Время на валидацию & Пресеты для распространённых задач \\
Нет официальных бенчмарков & Сложно сравнить & Внутренние тесты на LJSpeech \\
\bottomrule
\end{tabular}
\end{table}

\subsection{Направления будущей работы}

\begin{enumerate}
    \item \textbf{Автоматическая настройка гиперпараметров} — Bayesian Optimization для $\eta, \lambda_{decay}, \tau_{base}$.
    
    \item \textbf{Расширение на мультимодальность} — Интеграция визуального контекста.
    
    \item \textbf{Распределённый сон} — Консолидация на сервере для Edge-устройств.
    
    \item \textbf{Формальная верификация} — Доказательство устойчивости через Lyapunov-функции.
\end{enumerate}

\section{Лицензирование и воспроизводимость}

\textbf{Рекомендация:} Для обеспечения воспроизводимости исследований:

\begin{enumerate}
    \item Фиксировать версию кода (Git commit hash) для каждого эксперимента.
    \item Сохранять конфигурационные файлы вместе с весами модели.
    \item Публиковать seed для генераторов случайных чисел.
    \item Документировать версию зависимостей (torch, numpy, cuda).
\end{enumerate}

\textbf{Лицензия:} Рекомендуется открытая лицензия (MIT/Apache 2.0) для исследовательских целей с указанием авторства.

\section{Заключение раздела}

Данный раздел завершает техническую спецификацию DREAM-Net v1.0. Представленные интерфейсы, гиперпараметры и рекомендации по реализации предназначены для обеспечения корректной и воспроизводимой реализации архитектуры.

Все компоненты спецификации являются предметом дальнейшей экспериментальной валидации. Разработчикам рекомендуется начинать с базовой конфигурации и проводить абляционные тесты для настройки под конкретные задачи.
